\chapter{何时停止依赖AI：边界的识别}

在 Prompting 阶段，人通过任务描述和约束引导 AI；在 Modeling 阶段，人与 AI 共同完成数据处理、建模和实验；进入 Evaluation 阶段后，最重要的问题不再是“AI 还能继续做什么”，而是“现有证据允许我们相信到什么程度”。

AI 能够继续生成代码、图表和解释，并不代表它应当继续主导任务。当数据不足、知识缺失、统计结果不稳定，或者问题从预测跨越到因果与决策时，继续要求 AI 给出确定答案，往往只会得到更加完整但并不更加可靠的输出。

因此，识别边界是一种评价能力：判断 AI 当前适合担任执行者、副驾驶还是参谋；判断什么结论仍然有证据支持；以及在什么情况下必须停止、补充信息或交还给人判断。

% ----------------------------------------------------------------
\section{三类问题：AI擅长、AI可辅助、AI不应主导}
% ----------------------------------------------------------------

\subsection{AI作为执行者}

当任务定义明确、规则已经确认、结果容易验证时，AI 可以作为执行者。例如：

\begin{itemize}
  \item 按已确认规则生成数据质量检查代码；
  \item 实现指定的基准模型和评价指标；
  \item 绘制已经确定的验证图形；
  \item 整理实验结果并生成表格；
  \item 根据测试失败信息修改局部代码。
\end{itemize}

在这类任务中，人的主要工作是定义标准并检查结果。AI 的优势是快速实现，而不是替人决定什么算正确。

\subsection{AI作为副驾驶}

当任务存在多个合理方案，需要结合数据和领域知识进行选择时，AI 更适合作为副驾驶。它可以提出候选方案、解释差异、实现实验并总结结果，但关键选择仍由人确认。

常见任务包括：

\begin{itemize}
  \item 选择缺失值处理方式；
  \item 设计特征和模型候选集合；
  \item 选择数据划分和评价指标；
  \item 分析错误案例并提出改进方向；
  \item 判断复杂模型是否值得采用。
\end{itemize}

副驾驶模式的重点是把“AI 直接决定”改为“AI 生成候选，人根据证据选择”。

\subsection{AI作为参谋}

当问题依赖 AI 无法获得的信息，或者错误后果较高时，AI 只能作为参谋。它可以整理已有证据、提出需要调查的问题和比较不同观点，但不应主导最终判断。

例如：

\begin{itemize}
  \item 判断模型发现的相关关系是否代表因果关系；
  \item 决定高风险预测是否应当用于真实行动；
  \item 在证据冲突时选择应相信哪一方；
  \item 决定涉及安全、权利或重大资源分配的措施；
  \item 判断项目结果是否足以支持正式政策结论。
\end{itemize}

\subsection{角色应按步骤分配}

同一个项目中，AI 的角色会随任务变化。AI 可以作为执行者生成绘图代码，作为副驾驶分析模型错误，却只能作为参谋讨论是否应根据模型调整现实政策。

\begin{table}[htbp]
\centering
\small
\begin{tabular}{p{0.15\textwidth}p{0.24\textwidth}p{0.23\textwidth}p{0.18\textwidth}}
\toprule
角色 & 适用条件 & AI主要工作 & 人的主要工作 \\
\midrule
执行者 & 标准明确、容易验证 & 按规则实现与运行 & 定义标准、检查输出 \\
副驾驶 & 存在多个合理方案 & 生成候选、比较实验 & 选择方案、确认假设 \\
参谋 & 信息不足或后果较高 & 整理证据、提出问题 & 补充知识、承担决定 \\
\bottomrule
\end{tabular}
\end{table}

面对新任务，可以先使用下面的 Prompt 要求 AI 判断自身角色：

\begin{quote}
请判断这个任务中哪些步骤具有明确正确标准，哪些步骤需要领域判断，哪些步骤现有证据不足以作出结论。分别说明你适合担任执行者、副驾驶还是参谋，以及人需要确认什么。
\end{quote}

% ----------------------------------------------------------------
\section{AI的边界在哪里}
% ----------------------------------------------------------------

\subsection{知识边界：AI不知道项目中没有提供的信息}

AI 可能知道一般机器学习方法，却不知道本地数据如何采集、字段在组织内部如何定义、某次异常事件为何发生，也不知道没有被提供的现场经验。

例如，一个字段名为 \texttt{status=0}，AI 可能把它解释为“正常”，但在真实系统中它也可能表示“设备离线”。如果字段含义会改变分析结论，就不能让 AI 根据命名自行猜测。

遇到知识边界时，应补充数据字典、业务规则、现场说明或专业人员判断。无法补充时，应明确标记未知，而不是要求 AI 给出更确定的解释。

\subsection{数据边界：模型不能从不存在的数据中学习}

模型的能力受到训练数据覆盖范围限制。没有极端天气数据，就难以评价极端天气下的表现；没有新区域样本，就不能假设模型能够迁移到新区域；标签本身存在系统偏差时，模型可能只会复制这种偏差。

数据边界常见表现包括：

\begin{itemize}
  \item 关键类别样本过少；
  \item 训练数据没有覆盖真实使用场景；
  \item 标签定义不一致或质量不足；
  \item 重要变量从未被记录；
  \item 数据来自方便获取的对象，而非目标总体。
\end{itemize}

更多 Prompt 或更复杂模型不能自动补足缺失数据。此时应补充数据、缩小任务范围，或降低结论强度。

\subsection{统计边界：高分数不等于稳定规律}

模型在某次划分或某个指标上表现良好，不代表它发现了稳定规律。样本量较小、类别不平衡、重复尝试大量模型或只报告最好结果，都可能使性能看起来比实际更好。

当不同数据划分、不同时间窗口或不同指标给出明显不同的评价时，应承认统计证据仍不稳定。继续让 AI 调参可能得到更高分数，却不一定提高可信度。

\subsection{因果边界：预测关系不等于干预效果}

机器学习模型擅长利用相关关系进行预测，但“什么因素能够预测结果”与“改变什么因素能够改变结果”是两个不同问题。

例如，模型发现拥堵路段更容易发生事故，可以用于识别高风险时段；但不能仅凭这个结果断言减少流量一定会按预测幅度降低事故。现实干预还可能改变驾驶行为、路线选择和其他条件。

当任务从预测跨越到“为什么”或“采取什么措施会有效”时，应停止把模型输出直接解释为行动依据，并引入更合适的研究设计与领域判断。

\subsection{工程边界：模型正确不代表系统可用}

离线实验中的模型只处理准备好的数据。真实系统还会遇到字段缺失、输入延迟、软件版本变化、计算资源不足和上游服务中断。

一个模型即使预测性能良好，如果无法识别异常输入、无法在规定时间内运行，或者发生错误后没有基准方案，也不适合直接使用。工程边界要求评价者把“模型分数”扩展为“整个使用过程是否可靠”。

\subsection{责任边界：能够生成建议不等于能够承担决定}

AI 不承担现实决定产生的责任。涉及安全、隐私、权利、公共资源和重大成本时，不能因为 AI 给出了清晰建议，就把最终判断交给它。

责任边界并不意味着不能使用 AI，而是要求明确：谁理解证据，谁批准采用，谁监控结果，谁在出现问题时停止使用。

% ----------------------------------------------------------------
\section{如何识别AI已经到达边界}
% ----------------------------------------------------------------

\subsection{第一步：识别当前角色}

先判断当前步骤是在执行明确规则、比较候选方案，还是作出需要领域知识和责任承担的决定。如果任务已经从实现代码转向判断现实含义，AI 的角色通常需要降低。

\subsection{第二步：检查证据链是否完整}

一个可靠结论应能够回答：输入数据来自哪里，经过哪些处理，使用什么评价方法，结果如何支持结论，结论又适用于哪些条件。

可以要求 AI 生成证据表：

\begin{quote}
请把当前结论拆成若干条主张。对每条主张列出支持它的数据、代码输出、评价结果和假设，并指出哪些部分仍缺少证据。
\end{quote}

如果关键结论只能由 AI 的语言解释支持，却无法追溯到数据、实验或原始来源，就已经到达边界。

\subsection{第三步：检查所需数据和信息是否存在}

应明确哪些信息已经提供，哪些信息只是 AI 的推测。特别要检查：

\begin{itemize}
  \item 数据是否覆盖目标场景；
  \item 特征在预测时是否真实可用；
  \item 字段和标签含义是否得到确认；
  \item 是否存在无法观测但会改变结论的重要因素；
  \item 最新现场情况是否与历史数据一致。
\end{itemize}

缺少决定方向的信息时，正确动作是停止并请求补充，而不是采用默认假设。

\subsection{第四步：检查总体指标是否掩盖风险}

总体平均指标可能掩盖少数类别、关键区域或极端情景中的失败。应回到上一章的验证方法，检查基准、关键分组、错误案例和极端情景。

如果模型只在不重要的场景中优于基准，却在真正关心的场景中表现不足，就不能根据总体分数继续依赖它。

\subsection{第五步：检查结论的适用范围}

每个模型结论都应附带条件：适用于什么时间、区域、人群、数据质量和任务目标。如果准备把结果用于未经验证的新场景，就已经超出当前证据范围。

常见信号包括：

\begin{itemize}
  \item 使用对象与训练样本明显不同；
  \item 输入值超出训练数据范围；
  \item 数据采集方式或标签规则已经改变；
  \item 原本用于探索的模型被要求自动决策；
  \item 原本用于预测的结果被解释为因果结论。
\end{itemize}

\subsection{第六步：检查是否涉及因果、责任与高风险决定}

当任务涉及“应不应该采取措施”“谁应获得资源”“是否允许自动执行”等问题时，评价不能只依赖模型性能。此时需要补充现实约束、影响分析、专业判断和明确责任人。

如果这些条件尚未具备，就应停止 AI 主导，将其降为提供候选信息的参谋。

% ----------------------------------------------------------------
\section{如何处理边界：用领域知识重新组织协作}
% ----------------------------------------------------------------

发现边界后，不一定要放弃项目。更常见的做法是缩小问题、补充证据并重新安排人机分工。

\subsection{补充上下文，而不是让AI猜测}

当 AI 缺少字段定义、采集过程、业务规则或现场情况时，应先补充这些信息。可以要求 AI 列出它完成判断所需的最小信息集合，再由人逐项提供或标记无法获得。

\subsection{把领域知识写成明确约束}

专业判断如果只存在于人的经验中，AI 很难稳定使用。应尽量把它转化为规则，例如：

\begin{itemize}
  \item 哪些字段在预测时不可用；
  \item 哪些数值范围明显不合理；
  \item 哪些样本不能自动删除；
  \item 哪些关键场景必须单独评价；
  \item 哪些条件出现时必须停止输出。
\end{itemize}

这些约束既可以写入 Prompt，也可以实现为自动检查。

\subsection{补充证据或缩小结论}

如果数据不足以支持原问题，可以采取两种方向：收集更多数据，或者把结论缩小到现有证据能够支持的范围。

例如，模型没有在节假日数据上验证，就不应声称适用于全年；可以先将结论改为“适用于常规工作日”，并把节假日评价列为后续工作。

\subsection{降低AI角色}

当任务风险提高或证据变弱时，应把 AI 从执行者降为副驾驶，或从副驾驶降为参谋。例如：

\begin{itemize}
  \item 从“自动删除异常记录”改为“标记疑似异常记录供人检查”；
  \item 从“自动选择最终模型”改为“比较候选模型并呈现证据”；
  \item 从“自动执行控制策略”改为“生成建议供操作人员确认”；
  \item 从“给出政策结论”改为“列出可能解释和所需证据”。
\end{itemize}

\subsection{把决定改写为候选生成}

AI 最适合扩展候选空间，而不一定适合替人完成最后选择。面对不确定问题，可以要求它生成多个解释、多个方案和多个反例，并说明各自依赖的假设。人再结合领域知识、成本和责任进行选择。

\subsection{设置停止条件和分阶段确认}

复杂任务应在开始前定义停止条件。例如：

\begin{quote}
如果发现特征可能包含未来信息、模型未稳定优于基准、关键场景样本不足，或结论需要因果解释，请停止继续优化，列出当前证据和缺失信息，等待人工确认。
\end{quote}

还可以把任务拆成“理解数据、确认方案、实现模型、验证结果、解释结论”几个阶段，在方向容易修改时及时确认。

% ----------------------------------------------------------------
\section{案例：是否让AI自动调整信号配时}
% ----------------------------------------------------------------

某项目使用 AI 辅助建立短时交通流预测模型，希望根据预测结果自动调整路口信号配时。模型在历史测试集上的总体误差明显低于基准，单看分数似乎已经可以投入使用。

首先检查数据边界。训练数据主要来自常规工作日，没有覆盖大型活动和严重事故；部分路口传感器缺失较多。因此，模型在这些场景中的可靠性未知。

接着检查统计与工程边界。分组评价显示，模型在平峰时段表现较好，但在高峰时段改善有限。传感器中断时，模型仍会输出看似正常的预测，却无法提示输入异常。

然后检查因果边界。模型能够预测流量，却不能直接证明某种配时调整会改善整个路网。调整一个路口还可能把排队转移到相邻路口。

最后检查责任边界。自动调整会影响真实交通运行，错误操作的后果不能由模型承担。

因此，项目没有直接停止使用模型，而是重新安排了角色：

\begin{enumerate}
  \item AI 继续作为执行者生成预测和验证图形；
  \item AI 作为副驾驶提出配时调整候选方案；
  \item 调度人员结合现场情况决定是否采用；
  \item 传感器缺失、极端天气和预测超出合理范围时，系统停止建议并回到原有方案；
  \item 记录每次建议、人工决定和实际结果，为后续评价积累证据。
\end{enumerate}

这个案例说明，边界识别不是简单回答“能不能使用 AI”，而是根据现有证据确定 AI 可以参与到哪一步，以及什么时候必须把决定交还给人。

% ----------------------------------------------------------------
\section{实战练习}
% ----------------------------------------------------------------

\subsection*{练习一：识别AI角色}

选择一个机器学习项目，将数据理解、数据处理、建模、评价和结论解释分别标记为执行者、副驾驶或参谋任务。说明每一步为什么需要这种角色安排。

\subsection*{练习二：绘制边界地图}

针对一个模型结果，分别检查知识、数据、统计、因果、工程和责任边界。为每项边界写出当前证据、缺失信息和处理方法。

\subsection*{练习三：把决定改写为候选生成}

选择一个原本要求 AI 直接作决定的 Prompt，将其改写为生成候选方案、列出假设、提供验证方法并等待人工选择的 Prompt。

\subsection*{练习四：设计停止条件}

为一个即将用于真实场景的模型设计至少五条停止条件。条件应覆盖数据异常、表现低于基准、适用范围变化、因果解释和高风险操作。
